\documentclass[11pt]{amsart}

\usepackage[T1]{fontenc}
\usepackage{lmodern}
\usepackage{microtype}
\usepackage{amsmath,amssymb,amsthm}
\usepackage{graphicx}
\usepackage[
  left=0.90in,
  right=0.90in,
  top=0.72in,
  bottom=0.72in
]{geometry}
\usepackage[hidelinks]{hyperref}

\hypersetup{
  pdftitle={
    A Counterexample to He and Huang's Runge--Kutta Necessity Conjecture
  }
}

\newtheorem{theorem}{Theorem}
\newtheorem{proposition}[theorem]{Proposition}

\title{A Counterexample to He and Huang's Runge--Kutta\\
Necessity Conjecture}
\date{}

\subjclass[2020]{65L05, 65L06}
\keywords{
  explicit Runge--Kutta method, order conditions,
  simplifying assumptions, counterexample
}

\begin{document}

\begin{abstract}
He and Huang conjectured that the sufficient \(Q\)- and \(D\)-condition
package in their Theorem~1.2 is also necessary for Runge--Kutta order.
We disprove the conjecture under the printed Hadamard condition
\[
QO(m): b\odot Q_m=\{0\}.
\]
For every internally consistent explicit tableau and every admissible
order \(p\ge 5\), this condition forces \(b_2c_2^2=0\). The embedded
formula in Verner's pair \(\mathrm{RK}(8\!-\!6{:}5)a\) has order
exactly five and satisfies
\[
\widehat b_2=-\frac15,
\qquad
c_2=\frac15,
\]
so the proposed necessary package fails for every admissible pair
\((m,n)\). The argument concerns the printed Hadamard product and does
not apply to scalar orthogonality.
\end{abstract}

\maketitle

\section{The obstruction}

He and Huang define, for nonnegative integers \(r\),
\[
q_r
 =A c^{\odot r}
  -\frac{1}{r+1}c^{\odot(r+1)},
\]
and
\[
Q_1=\operatorname{span}\{q_0\},
\qquad
Q_2=\operatorname{span}
 \{q_1,Aq_0,q_0\odot c,q_0\}.
\]
Their spaces satisfy
\[
Q_1\subseteq Q_2\subseteq Q_3\subseteq\cdots
\]
\cite[Defs.~3.4 and 3.9, Eq.~(16)]{HeHuangTheory}.
The two entries in the bibliography below share the arXiv identifier
2605.16995: that record was retitled and reduced to a technical appendix
at v3, and v3 carries no definitions or notation of its own, recalling
the general theory instead from the v1-titled work, which it cites under
that title. We therefore cite definitions and notation to v1, and
Theorem~1.2 and Conjecture~2.1 to v3.

Theorem~1.2 of arXiv:2605.16995v3 states that an explicit Runge--Kutta
method has order at least \(p\) if, for some nonnegative integers \(m,n\)
with
\[
m\ge n-1,
\qquad
m+n+1\ge p,
\]
six conditions hold: the quadrature conditions \(B(p)\), the
\(Q\)-orthogonality conditions
\[
QO(m):
\qquad
b\odot Q_m=\{0\},
\]
the \(D\)-orthogonality conditions \(DO(n)\), the weak mutual
orthogonality conditions \(QD_{\mathrm{weak}}(m,n)\), the pivot residual
condition \(PR(n)\), and the \(Q\)-ring condition \(QR(m)\)
\cite[Thm.~1.2]{HeHuangDetails}. Conjecture~2.1 reads, verbatim
\cite[Conj.~2.1]{HeHuangDetails}:
\begin{quote}
The set of conditions in [Theorem~1.2] is not only sufficent [sic] but also
necessary for a vector-valued explicit/implicit RK method to achieve order
\(p\).
\end{quote}
Here \(\odot\) is the componentwise Hadamard product; the same source
uses \(\cdot\) for the scalar inner product
\cite[Notations~2.1 and 2.3]{HeHuangTheory}.
The Hadamard reading of \(QO(m)\) is forced by the source rather than
inferred: all set operations there are declared to be Minkowski
operations, so \(b\odot Q_m=\{0\}\) abbreviates
\(\{b\odot q:q\in Q_m\}=\{0\}\), that is, \(b\odot q=0\) as a vector for
every \(q\in Q_m\) \cite[Notation~2.5]{HeHuangTheory}. The choice of
product is deliberate within Theorem~1.2 itself, whose conditions
\(DO(n)\) and \(QD_{\mathrm{weak}}(m,n)\) are written with \(\cdot\) and
not with \(\odot\).

Necessity of the package asserts that every one of the six conditions
holds whenever the method has order \(p\), so the conjecture is refuted
by a single method of order \(p\) that violates a single condition for
every admissible \((m,n)\). We exhibit one violating \(QO(m)\).

\begin{proposition}\label{prop:obstruction}
Let \((A,b,c)\) be an internally consistent explicit Runge--Kutta
tableau, so that \(A\) is strictly lower triangular and
\(c=A\mathbf 1\). Let \(p\ge 5\). If nonnegative integers \(m,n\)
satisfy
\[
m\ge n-1,
\qquad
m+n+1\ge p,
\]
and \(QO(m)\) holds, then
\[
b_2c_2^2=0.
\]
\end{proposition}

\begin{proof}
Since \(A\) is strictly lower triangular its first row vanishes, so internal
consistency gives \(c_1=(A\mathbf 1)_1=0\). The vector \(q_1\) is one of the
listed generators of \(Q_2\), so in particular \(q_1\in Q_2\). The second
row of \(A\) has \(a_{21}\) as its only possibly nonzero entry, so
\[
(q_1)_2
 =(Ac)_2-\frac12c_2^2
 =a_{21}c_1-\frac12c_2^2
 =-\frac12c_2^2.
\]

Moreover, \(m\ge n-1\) implies \(n\le m+1\), whence
\[
p\le m+n+1\le 2m+2.
\]
Because \(p\ge 5\), one has \(m\ge 2\). Thus
\[
q_1\in Q_2\subseteq Q_m,
\]
and \(QO(m)\) gives
\[
0=(b\odot q_1)_2
  =-\frac12b_2c_2^2.
\]
\end{proof}

\section{An order-five counterexample}

\begin{theorem}\label{thm:counterexample}
Under the printed Hadamard definition of \(QO(m)\), Conjecture~2.1 in
arXiv:2605.16995v3 is false. It already fails for an explicit
Runge--Kutta method of order exactly five.
\end{theorem}

\begin{proof}
Consider the embedded formula in Verner's eight-stage pair
\(\mathrm{RK}(8\!-\!6{:}5)a\)
\cite[Table~4 and Algorithm~1.2]{Verner2014}:
\begin{center}
\begingroup
\renewcommand{\arraystretch}{1.35}
\resizebox{0.99\textwidth}{!}{%
\(
\begin{array}{c|cccccccc}
0\\
\frac15
 & \frac15\\
\frac14
 & \frac{3}{32}
 & \frac{5}{32}\\
\frac47
 & \frac{60}{343}
 & -\frac{440}{343}
 & \frac{576}{343}\\
\frac79
 & \frac{847}{17496}
 & \frac{770}{729}
 & -\frac{17024}{19683}
 & \frac{84721}{157464}\\
\frac15
 & \frac{523}{2240}
 & 0
 & -\frac5{57}
 & \frac{245}{2496}
 & -\frac{1215}{27664}\\
1
 & \frac{15}{2212}
 & \frac{314}{395}
 & \frac{3448}{1501}
 & -\frac{2695}{4108}
 & \frac{203391}{273182}
 & -\frac{864}{395}\\
\frac47
 & \frac{1185}{10976}
 & -\frac{2543}{1715}
 & \frac{36262}{19551}
 & -\frac{245}{1248}
 & \frac{1215}{13832}
 & \frac15
 & 0\\
\hline
 & \frac{43}{560}
 & -\frac15
 & \frac{2816}{7695}
 & -\frac{41}{84240}
 & \frac{19683}{69160}
 & \frac15
 & \frac{79}{1080}
 & \frac15
\end{array}
\)
}%
\endgroup
\end{center}

Let \(\widehat b\) denote the displayed weight vector. The order
conditions are indexed by rooted trees: with elementary weights
\(\Phi(\bullet)=\mathbf 1\) and
\(\Phi([t_1,\dots,t_k])_i=\prod_{j=1}^k(A\Phi(t_j))_i\), and densities
\(\gamma(\bullet)=1\) and \(\gamma(t)=|t|\prod_{j=1}^k\gamma(t_j)\), the
method has order \(p\) precisely when
\(\widehat b^{\mathsf T}\Phi(t)=1/\gamma(t)\) for every rooted tree \(t\)
with \(|t|\le p\). There are \(1,1,2,4,9\) such trees at orders one
through five, and all seventeen conditions hold exactly on the displayed
tableau; of the twenty trees of order six, eight give conditions that
fail. For instance, the chain conditions of orders three, four and five
hold,
\[
\widehat b^{\mathsf T}Ac=\frac16,
\qquad
\widehat b^{\mathsf T}A^2c=\frac1{24},
\qquad
\widehat b^{\mathsf T}A^3c=\frac1{120},
\]
while the one of order six fails:
\[
\widehat b^{\mathsf T}A^4c
 =\frac{1}{672}
 \ne\frac{1}{720}.
\]
Thus the formula has order exactly five.

For the displayed tableau,
\[
Ac
 =\left(
  0,0,\frac1{32},\frac8{49},\frac{49}{162},
  0,\frac12,\frac8{49}
 \right)^{\mathsf T},
\]
so that \((Ac)_i=\frac12c_i^2\) at the six stages \(i=1,3,4,5,7,8\),
whereas the two remaining stages both have \(c_i=\frac15\) and
\((Ac)_i=0\). Hence
\[
q_1
 =Ac-\frac12c^{\odot2}
 =\left(
  0,-\frac1{50},0,0,0,-\frac1{50},0,0
 \right)^{\mathsf T}.
\]
Consequently,
\[
\widehat b\odot q_1
 =\left(
  0,\frac1{250},0,0,0,-\frac1{250},0,0
 \right)^{\mathsf T}
 \ne0.
\]

At \(p=5\), every admissible pair \((m,n)\) has \(m\ge 2\), so
\(q_1\in Q_m\) and \(QO(m)\) fails for every such pair. The method has
order five nevertheless, disproving the conjectured necessity.
\end{proof}

\paragraph{Scope.}
For the same formula,
\[
\widehat b^{\mathsf T}q_1=0.
\]
More generally, the third-order conditions imply
\[
b^{\mathsf T}q_1
 =b^{\mathsf T}Ac-\frac12b^{\mathsf T}c^{\odot2}
 =\frac16-\frac12\cdot\frac13
 =0.
\]
Thus the result concerns the coordinatewise condition
\(b\odot Q_m=\{0\}\) as printed and does not refute a modified
conjecture based on scalar orthogonality. The sufficiency theorem
itself is unaffected.

\begin{thebibliography}{9}

\bibitem{HeHuangTheory}
J.~He and J.~Huang,
\emph{Low stage high order explicit Runge--Kutta methods via
\(Q\)- and \(D\)-conditions: general theory and efficient recursive
construction},
arXiv:2605.16995v1, 2026.

\bibitem{HeHuangDetails}
J.~He and J.~Huang,
\emph{Low stage and high order explicit Runge--Kutta methods via
\(Q\)- and \(D\)-conditions: several construction details},
arXiv:2605.16995v3, 2026.

\bibitem{Verner2014}
J.~H.~Verner,
\emph{Explicit Runge--Kutta pairs with lower stage-order},
Numer. Algorithms \textbf{65} (2014), no.~3, 555--577,
\href{https://doi.org/10.1007/s11075-013-9783-y}
{doi:10.1007/s11075-013-9783-y}.

\end{thebibliography}

\end{document}